\section{Discussion}
\label{sec:discussion}

\subsection{Answering the research questions}

\paragraph{RQ1 (BOND acceleration): answered in the negative for this
setting.} Exact-safe BOND dimension pruning did not accelerate exact ColBERT
MaxSim over the project's compiled exhaustive kernel. On raw 128-dimensional
embeddings the kernel retained 95.47\% (SciFact) and 94.16\% (NFCorpus) of the
exhaustive component products and ran 10.76$\times$ and 9.84$\times$ slower
than the precision-matched exhaustive arm of the same interleaved run. The
result is a controlled measurement with identical inputs, precision, thread
budgets, and timing boundaries in both arms, and it replicates across two
datasets.

\paragraph{RQ2 (bound behavior).} The residual Cauchy--Schwarz bound stays
loose because normalized ColBERT embeddings spread score energy almost
uniformly across dimensions: on raw embeddings, every SciFact prune and almost
every NFCorpus prune occurred at checkpoint 96 of 128, after 75\% of the
arithmetic was already spent. The PCA experiment isolates the two remaining
obstacles. Concentrating 90.91\% of energy in eight components moves pruning
earlier and removes 31--38\% of products, so dimension order is genuinely part
of the problem; but the persisting 8.38--8.51$\times$ slowdown shows that
bound-maintenance and memory-irregularity costs, not multiply--accumulate
volume, dominate this kernel. An exact-safe pruning kernel for MaxSim would
need per-pair bookkeeping that is nearly free before it could profit from
even a favorable dimension order.

\paragraph{RQ3 (flat token index as candidate generator): answered in the
affirmative, with a dataset-dependent ceiling.} With $L = 100$ and
$\texttt{nprobe} = 8$, the flattened token index produced candidate pools
containing the complete exact top-10 for all 40 SciFact queries and for 37 of
40 NFCorpus queries. Token-hit aggregation is not a valid final scorer
(\cref{sec:three-stages}), but as a selector feeding exact MaxSim reranking
it supports a smooth, controllable quality--work curve.

\paragraph{RQ4 (PDX-IVF versus FAISS-IVF).} The two engines produced
identical candidate pools, identical exact-ranking recovery curves, and no
consistent latency ordering across budgets and datasets; PDX token search was
slightly slower in every matched NFCorpus arm while total latencies were
nearly tied. The evidence indicates that, in this workload and at these corpus
sizes, the vertically decomposed IVF implementation neither helps nor hurts
candidate generation relative to a conventional IVF.

\paragraph{RQ5 (controlled PLAID comparison).} The two PLAID configurations
selected on 10 SciFact validation queries do not transfer: validation exact
recall@10 of 0.89/0.92 falls to 0.2875/0.4375 on held-out SciFact and to
0.3100/0.3950 on NFCorpus. They are also slower than compiled exact MaxSim on
these small CPU workloads. The post-hoc full-score sensitivity arms improve
recovery to 0.6750 and 0.5975, but remain approximate and take 24.65 and 35.86
times their same-run exact medians. Thus the tested PLAID configuration is
dominated here, but the evidence is deliberately limited to the pinned
four-bit CPU implementation and corpus scale rather than generalized to
PLAID's intended large-scale setting.

\subsection{Interpretation and scope of claims}

Three boundaries on interpretation deserve emphasis. First, the IVF pipeline
is effective but \emph{standard}: candidate generation followed by exact
reranking is textbook approximate-retrieval architecture, and since FAISS
reproduces the identical curve, none of that effectiveness is attributable to
PDX. Second, no end-to-end speedup is claimed anywhere in this report. The
online boundary excludes query encoding and index construction, and the IVF
and BOND experiments use different precision contracts, so the only defensible
latency statements are within-run, precision-matched ratios. Third, the
measured negative BOND result is a statement about the tested exact-safe
kernel on the tested embeddings---a document-oriented research prototype, not
an upstream PDX block kernel---so it bounds this design point rather than the
entire idea of dimension pruning for late interaction.

PLAID adds a fourth boundary: its \texttt{n\_full\_scores} parameter is a
configured budget, while the current API hides realized candidates and
intermediate stages. The selected points are valid complete-boundary latency
measurements, but their work cannot be equated one-for-one with IVF's observed
rerank counts. The full-score points are post-hoc mechanism checks and are not
used to claim a tuned held-out frontier.

Within those boundaries, the practical guidance for multi-vector retrieval
systems is clear: engineering effort is better spent on candidate generation
quality (where NFCorpus shows a real coverage ceiling) and on fast exhaustive
reranking kernels than on exact-safe dimension pruning of raw ColBERT
representations.
